\section{Что установлено, какие кандидаты исключены и что остаётся открытым}
\label{sec:status-literature}

\paragraph{Статус исследовательской ветки.}
Все утверждения этого приложения представлены как рабочие доказательства и
вычислительные сертификаты, находящиеся на стадии проверки. Внутренняя
проверка тождеств и повторное применение опубликованных критериев не равны
внешнему независимому рецензированию. Ниже слово \emph{следствие} обозначает
логический вывод из выписанных формул и указанных теорем; оно не означает
заявления о литературном приоритете.

\begin{center}
\small
\begin{tabular}{@{}p{0.30\linewidth}p{0.64\linewidth}@{}}
\toprule
Вопрос & Результат и его точная граница\\
\midrule
Конкретная кривая Шастри &
Полиномиальный выпрямитель над $K=\Q(\sqrt{-3})$ исключает эту
параметризацию из кандидатов на комплексную невыпрямляемость. Он не
исключает существования других невыпрямляемых пространственных прямых.\\[4pt]
Гармоническая модель и $C_N$ &
Выпрямляются именно указанные параметризации. Гармоническая модель
использует $D_5=t^5-5t^3+5t$; замена $5t$ на $4t$ не покрывается
этим утверждением.\\[4pt]
Ручность построенных карт &
Дикость $\Phi_T,\Phi_H,\Phi_N$ следует из критериев относительных
координат. Отсутствие \emph{любого} ручного выпрямителя тех же кривых
не доказано.\\[4pt]
Стабилизация &
$\Phi_\ast\times\id_{\A^3}\in\Tame_6(K)$ следует из известной
теоремы о стабильной ручности. Минимальное число дополнительных
переменных не найдено.\\[4pt]
Кремонова подгруппа &
Построенные карты принадлежат $\mathcal C_3(K)$.
Это не описание всей $\Cr_3(K)$ и не первый пример дикого элемента
в $\mathcal C_3(K)$.\\[4pt]
Поле определения &
Для Шастри минимальная степень подходящего числового поля над $\Q$
равна $2$. Единственность подходящего квадратичного поля и эффективная
классификация полей выпрямления не установлены.\\
\bottomrule
\end{tabular}
\end{center}

\paragraph{Исторический кандидат и общая проблема.}
В работе Шастри \cite[\S1 и \S3, Example~1]{Shastri1992} вещественные
узловые параметризации предложены как кандидаты для исследования
комплексной невыпрямляемости. Поэтому наличие полного полиномиального
автоморфизма для $\gamma_T$ решает вопрос о \emph{данной кривой}
положительно. Историческая гипотеза о существовании хотя бы одного
невыпрямляемого вложения $\A^1_{\C}\hookrightarrow\A^3_{\C}$ этим не
опровергается. Равным образом не доказана противоположная универсальная
гипотеза о выпрямляемости всех таких вложений. Мы не приписываем статье
Шастри отдельную нумерованную гипотезу, которой там не сформулировано.

\paragraph{Связанные гипотезы из первого приложения.}
Нужно отдельно учитывать уже записанные в первом приложении следствия
для кандидатов ван ден Эссена--ван Россума
\cite[\S5, Conjectures~5.1--5.3]{VdER2002}.\footnote{Нумерация и
сопоставление объектов здесь унаследованы из Appendix 1. При подготовке
этой версии издательская запись была проверена, но повторная загрузка
полного текста AMS оказалась недоступна. Алгебраические выводы ниже
проверены непосредственно; это ограничение библиографического аудита
зафиксировано также в \texttt{audit/claim\_status.md}.}
Это продолжение исходного выпрямления, а не новый результат о группе
Кремоны. При $\gamma_T=(f,g,h)$ рассматривается
\[
 D=t\partial_u+f'(u)\partial_x+g'(u)\partial_y+h'(u)\partial_z
 \quad\text{на }\C[t,u,x,y,z].
\]
Гипотеза~5.1 (в последующей литературе также \emph{Conjecture E})
предсказывала неполиномиальность $\ker D$. Из выпрямителя
$\Phi_T=(P,Q,R)$ с обратным $\Psi$ получаются, при
$V=\gamma_T(u)-t(x,y,z)$,
\[
 p=P(V),\qquad q=Q(V)/t,\qquad r=R(V)/t,\qquad
 s=(u-P(V))/t.
\]
Все деления точны, поскольку $\Phi_T\gamma_T=(u,0,0)$. Обратная замена
равна
\[
 u=p+ts,\qquad
 (x,y,z)=\frac{\gamma_T(p+ts)-\Psi(p,tq,tr)}{t};
\]
она также полиномиальна. Следовательно,
\[
 \C[t,u,x,y,z]=\C[t,p,q,r,s],\qquad
 D=\partial_s,\qquad \ker D=\C[t,p,q,r].
\]
Это даёт отрицание именно указанного кандидата~5.1. Если $\sigma$
обозначает исходный слайс из того же построения, то
$\sigma-s\in\ker D$; отсюда
$\C[t,u,x,y,z]/(\sigma)\simeq\C^{[4]}$.
Также при $I_{\gamma_T}=(Q,R)$ имеем
\[
 \C[x,y,z][t,I_{\gamma_T}/t]
 =\C[t,P,Q/t,R/t]\simeq\C^{[4]}.
\]
Это соответствующие полиномиальные описания объектов~5.2 и~5.3;
подробные формулы, сопоставление обозначений и сертификаты находятся
в первом приложении репозитория. Общие проблемы сокращения,
вложений и линеаризации действий из этих частных вычислений не решаются.
Отозванный препринт Оды \cite{OdaWithdrawn} не используется как
доказательство: его запись arXiv содержит отзыв от 9 мая 2020 года
с указанием ошибки.

\paragraph{Предшественники программы.}
Связь задачи вложений с кремоновой эквивалентностью явно обсуждается
у Меллы--Поластри \cite[введение и Theorem~1]{MP2009}.
Относительные координаты, факторизации над полем частных и перенос
изучаются у Эдо--Венеро \cite{EV2001}.
Наиболее близкий предшественник конструктивной части --- Льюис
\cite[Lemma~1, Theorems~2 и~12]{Lewis2024}: координата, ограничивающаяся
в параметр, строится с использованием остаточных координат и
исправления рациональных преобразований. В этой же работе сопоставлены
предшествующие результаты Крайгеро, Бхатвадекара--Роя и Куроды.
Поэтому общий приём ``локализация, затем устранение знаменателей''
не заявляется здесь как новая стратегия.

Наличие диких автоморфизмов в трёх переменных давно установлено;
перенос дикости относительной координаты в три переменные даёт
\cite{UY2004}. Блан--Хеден уже доказали
$\operatorname{Nagata}\in\mathcal C_3$
\cite[Proposition~6.8]{BH2015}. Стабильная ручность используемого
класса следует из \cite[Corollary~4.9]{BEW2012}.
Самостоятельное содержание данного приложения состоит в применении
этих результатов к \emph{конкретным} выпрямителям, проверке их
гипотез и явном описании дальнейшей программы.

\paragraph{Нерешённые вопросы и границы поиска.}
Открытыми в нашей работе остаются полнота относительного метода,
существование ручного выпрямителя для тех же кривых, минимальные степени
выпрямителей, спуск в другие вполне мнимые числовые поля и общий случай
$\A^1\hookrightarrow\A^3$ над $\C$.
Конечные вычисления для $N\leq N_{\max}$ служат проверкой реализации;
утверждения для всех $N$ опираются на индуктивные доказательства.
Неуспех ограниченного поиска не доказывает невыпрямляемость.
Целевой литературный поиск по состоянию на 28 октября 2026 года не
выявил более ранней публикации с нашим полным выпрямителем Шастри,
однако отсутствие найденной работы не устанавливает приоритет.
